\section{Exact Classification and Residual Certificates}\label{app:exact-certificates}

\subsection{Period-eight orbit table}

Use \(0\) for negative flux and \(1\) for positive flux. The 128 legal period-eight
words form 18 dihedral orbits. The structural proof in Section~\ref{sec:eight-barrier} gives the
following exact classification.

\begin{table}[tbp]
\centering
\caption{Exact classification of legal period-eight flux orbits.}
\label{tab:period-eight-classification}
\begingroup
\renewcommand{\arraystretch}{1.12}
\scriptsize
\resizebox{\textwidth}{!}{%
\begin{tabular}{lllll}
\toprule
canonical \(Q\) & \(d\) & orbit size & primitive \(\tau\) period & exact conclusion \\
\midrule
\(00000000\) & 0 & 1 & 2 & \(R=8\) \\
\(00000011\) & 2 & 8 & 8 & \(R>8\) \\
\(00000101\) & 2 & 8 & 8 & \(R>8\) \\
\(00001001\) & 2 & 8 & 8 & \(R>8\) \\
\(00001111\) & 4 & 8 & 8 & \(R>8\) \\
\(00010001\) & 2 & 4 & 8 & \(R=\eta\) \\
\(00010111\) & 4 & 16 & 8 & \(R>8\) \\
\(00011011\) & 4 & 8 & 8 & \(R>8\) \\
\(00100111\) & 4 & 8 & 8 & \(R>8\) \\
\(00101011\) & 4 & 16 & 8 & \(R>8\) \\
\(00101101\) & 4 & 8 & 8 & \(R>8\) \\
\(00110011\) & 4 & 4 & 4 & \(R>8\) \\
\(00111111\) & 6 & 8 & 8 & \(R>8\) \\
\(01010101\) & 4 & 2 & 4 & \(R>8\) \\
\(01011111\) & 6 & 8 & 8 & \(R>8\) \\
\(01101111\) & 6 & 8 & 8 & \(R>8\) \\
\(01110111\) & 6 & 4 & 8 & \(R>8\) \\
\(11111111\) & 8 & 1 & 1 & \(R>8\) \\
\bottomrule
\end{tabular}
}
\endgroup
\end{table}

The orbit sizes sum to 128. The unique sub-eight orbit is \(00010001\), and the
unique equality orbit at eight is \(00000000\).

\subsection{Five low-period residual certificates}

For every row, reconstruct the canonical triangle lift by setting
\(\tau_{0}=+1\) and recursing through \(\tau_{i+1}=Q_i \tau_i\). The vector coordinates
are in the ordered fiber basis \((v_{0},\ldots,v_{p-1})\) of Section~\ref{subsec:floquet-matrices}. This fixes
the matrix denoted by \(H_Q(z)\) below without any gauge ambiguity.

For each row below, \(H=H_Q(z)\) and the listed integer vector is denoted by
\(v\). Direct matrix multiplication gives the stated rational quotient
\(r=\lVert Hv\rVert^{2}/\lVert v\rVert^{2}\).

\begin{table}[tbp]
\centering
\caption{Exact Rayleigh data for the five residual competitors.}
\label{tab:residual-rayleigh-data}
\begingroup
\renewcommand{\arraystretch}{1.12}
\begin{tabular}{llll}
\toprule
\(p\), \(Q\) & \(z\) & vector \(v\) & \(r\) \\
\midrule
10, \(0000010001\) & -1 & \((3,2,-3,3,0,5,-5,5,-5,2)\) & \(\frac{1066}{135}\) \\
12, \(000000010001\) & 1 & \((-5,-5,-5,-1,-4,-2,-3,0,-2,-2,0,-4)\) & \(\frac{1022}{129}\) \\
14, \(00000000010001\) & 1 & \((-1,2,0,2,-1,3,-2,3,-1,4,-4,4,-3,0)\) & \(\frac{119}{15}\) \\
14, \(00010010001001\) & 1 & \((-4,-5,-4,0,-2,0,1,0,-2,0,-4,-5,-4,-6)\) & \(\frac{1270}{159}\) \\
16, \(0000000000010001\) & 1 & \((-5,-5,-5,-1,-4,-2,-4,-1,-3,-1,-2,0,-2,-2,0,-4)\) & \(\frac{1204}{151}\) \\
\bottomrule
\end{tabular}
\endgroup
\end{table}

All five quotients exceed four. For the transformation

\begin{equation}
u=((r-4)^{2}-10)/2,
\end{equation}

the exact comparison data are:

\begin{table}[H]
\centering
\caption{Exact radical-comparison data for the residual competitors.}
\label{tab:residual-radical-data}
\begingroup
\renewcommand{\arraystretch}{1.35}
\begin{tabular}{lll}
\toprule
\(r\) & \(u\) & \(u^{2}-5\) \\
\midrule
\(\frac{1066}{135}\) & \(\frac{47213}{18225}\) & \(\frac{568314244}{332150625}\) \\
\(\frac{1022}{129}\) & \(\frac{44813}{16641}\) & \(\frac{623590564}{276922881}\) \\
\(\frac{119}{15}\) & \(\frac{1231}{450}\) & \(\frac{502861}{202500}\) \\
\(\frac{1270}{159}\) & \(\frac{74573}{25281}\) & \(\frac{2365487524}{639128961}\) \\
\(\frac{1204}{151}\) & \(\frac{65995}{22801}\) & \(\frac{1755912020}{519885601}\) \\
\bottomrule
\end{tabular}
\endgroup
\end{table}

Every entry in the last two columns is positive. Because \(r>4\), the branch
argument \eqref{eq:residual-radical-comparison} proves \(r>\eta\), and Rayleigh's principle proves \(R(Q)>\eta\).
These five computations are the only individual endpoint certificates needed
by the compressed proof of Theorem~\ref{thm:low-period-frontier}.

The other triangle lift is obtained from the canonical one by the conjugacy
and fiber reparametrization of Lemma~\ref{lem:operator-equivalences}; transporting \(v\) through that unitary
map gives the same quotient and conclusion.

\subsection{Exact two-defect moment table}

For a period-eight word with two positive fluxes at cyclic separation \(s\),
the dynamic program \eqref{eq:walk-recurrence}-\eqref{eq:walk-moment} gives the following exact moments and excesses.
The table includes every value needed to identify the first positive excess;
only the positive entries are used to infer \(R>8\).

\begin{table}[H]
\centering
\caption{Moment and excess sequences for separated period-eight defect pairs.}
\label{tab:separated-defect-moments}
\begingroup
\renewcommand{\arraystretch}{1.12}
\scriptsize
\resizebox{\textwidth}{!}{%
\begin{tabular}{lll}
\toprule
\(s\) & sequence & \(k=1,\ldots,10\) \\
\midrule
1 & \(M_k\) & \(32, 192, 1376, 10976, 93312, 823920, 7447136, 68348832, 633982976, 5926419872\) \\
1 & \(F_k\), \(k=1,\ldots,9\) & \(-64, -160, -32, 5504, 77424, 855776, 8771744, 87192320, 854556064\) \\
2 & \(M_k\) & \(32, 192, 1328, 9888, 76832, 612624, 4965328, 40683296, 335872832, 2788389472\) \\
2 & \(F_k\), \(k=1,\ldots,9\) & \(-64, -208, -736, -2272, -2032, 64336, 960672, 10406464, 101406816\) \\
3 & \(M_k\) & \(32, 192, 1280, 9056, 66592, 503088, 3877920, 30363808, 240761792, 1928966432\) \\
3 & \(F_k\), \(k=1,\ldots,9\) & \(-64, -256, -1184, -5856, -29648, -146784, -659552, -2148672, 2872096\) \\
4 & \(M_k\) & \(32, 192, 1280, 8928, 63872, 464496, 3417152, 25354656, 189356672, 1421345952\) \\
4 & \(F_k\), \(k=1,\ldots,9\) & \(-64, -256, -1312, -7552, -46480, -298816, -1982560, -13480576, -93507424\) \\
\bottomrule
\end{tabular}
}
\endgroup
\end{table}

Thus separations one, two, and three cross the barrier first at \(F_{4}\), \(F_{6}\),
and \(F_{9}\), respectively. The nonpositive displayed values for separation four
are not used in the converse direction; that case is settled by the exact
Floquet theorem of Section~\ref{sec:period-eight-edge}.

\subsection{The order-32 positive-definiteness cross-check}

For the witness \eqref{eq:order-thirty-two-triangle-word}, form \(M=1561I-200A_{32}^{2}\). The exact certificate consists
of the 32 positive pivots of rational \(LDL^T\) elimination. Fraction-free
Bareiss elimination independently produces the leading principal minors, and
the cumulative products of the \(LDL^T\) pivots agree with those minors. This
double route verifies both arithmetic consistency and the hypothesis of
Sylvester's criterion. The conclusion is the strict inequality
\(\rho(A_{32})^{2}<\frac{1561}{200}\); no approximation to an eigenvalue is part of the
certificate.

This elimination is an independent arithmetic cross-check, not the logical
proof used by Proposition~\ref{prop:order-thirty-two-witness}. The proposition follows directly from the
displayed fiber matrix \eqref{eq:target-floquet-matrix}, determinant \eqref{eq:floquet-polynomial}, positive expansion \eqref{eq:rational-positive-expansion},
and finite Floquet decomposition \eqref{eq:target-finite-decomposition}.
